\section{Discussion}

\subsection{Implications for astronomical software design}

Native Rust types connect Siderust's numerical components with explicit
astronomical semantics. Units, epochs, frames and centers participate in the
interfaces that join a calculation's stages.
Developers can therefore follow both the numerical data flow and the scientific
interpretation of its intermediate values. Explicit transformations provide
natural places to select a reference system, epoch or model context.

The observing applications illustrate this benefit in two settings. The
terrestrial workflow connects UTC input, TT and UT1 arguments, sidereal rotation
and altitude selection. The lunar workflow connects a kernel-derived position,
a selenocentric origin, body orientation and observer-relative projection.
Their shared representations support coherent application code while allowing
the astronomical models to differ.

\subsection{Genericity and extensibility}

Generic surface observers separate local geometry from the physical model of
the body. A body supplies its center, fixed frame and reference surface; the
common pipeline performs observer subtraction and horizontal projection.
The same structure serves a spherical Moon and an ellipsoidal Mars, including
both planetocentric and planetographic site representations. This is a concrete
example of extending astronomical capability by composing existing abstractions.

The modular structure also supports partial adoption. Applications may use
\texttt{qtty}, \texttt{affn} or \texttt{tempoch} directly, or combine numerical
components with their own models. A shared geometric vocabulary reduces
translation between component-specific containers, while explicit contexts
allow applications to control scientific data and model selection.

An existing downstream application is NSB v0.3.0, for night-sky radiance
evaluation and observation planning \cite{Valles_Puig_NSB_Night-Sky_Background}.
Its public observer is \texttt{Geodetic<ECEF>}; target directions retain their
frame, epochs use Tempoch's \texttt{Time<UTC>}, and radiometric quantities use
Siderust's quantity re-exports. NSB adds sky-model composition, radiance evaluation
and planning logic on this shared vocabulary. The tagged version depends on
Siderust v0.15.1, so it demonstrates existing ecosystem reuse rather than
validation of v0.16.0 or use of its new surface APIs. NSB shares authorship with
Siderust; it is not independent external adoption, and its observational results
are not used as a Siderust accuracy reference.

\subsection{Performance and practical integration}

The matched arithmetic measurements show that expressive dimensional interfaces
can compile to costs close to scalar arithmetic for the tested operations.
This result is useful for applications that perform unit-bearing calculations
frequently. Complete observing and dynamics workloads also include model
evaluation, context access and data loading, which motivate further workflow
benchmarks across hardware and application settings.

C-compatible interfaces extend the architecture's reuse to other languages by
sharing the Rust computation. Host-language adapters additionally handle
marshalling, resource ownership and error translation. Pinning wrapper releases
and testing the corresponding host-language workflows will extend the present
Rust evaluation into a
cross-language assessment.

\subsection{Compiler-guided development}

Typed interfaces also provide structured feedback during implementation.
Type-directed synthesis uses specifications to constrain program construction
\cite{polikarpova2016}; systems such as RustAssistant use compilation feedback
for iterative repair of Rust programs \cite{deligiannis2025rustassistant}.
Siderust supplies a scientific instance of this mechanism: incompatible units,
frames, centers or scales produce deterministic compiler diagnostics.

Our frame-mismatch fixture produces \texttt{E0277} when subtracting positions
in different frames. Listing~\ref{lst:frame-repair} repairs the operation with
an explicit transformation while retaining its frame, center and unit
information. These reproducible diagnostics can guide both a developer and
an automated programming tool through an iterative compile--revise--test loop.
A successful compilation confirms the expressed contracts; choosing a
physically suitable model and consistent epoch remains the developer's task.
Together with runtime checks and numerical validation, these interfaces provide a
structured foundation for constructing astronomical applications. Measuring
their effect on AI-assisted development is a further research direction.

\lstinputlisting[float=tbp,caption={Explicit frame alignment before position subtraction. Type aliases improve readability while preserving the scientific distinctions.},label={lst:frame-repair}]{../examples/src/bin/frame_repair.rs}

\subsection{Scientific scope and future directions}

The evaluation checks the specified operations and reference models over
deterministic samples. Same-kernel agreement validates the implemented
ephemeris path rather than the underlying physical ephemeris. Lunar pointing
uses analytical \texttt{IAU\_MOON}, with accuracy tied to that orientation
model and the chosen reference surface. Apparent-place corrections and terrain
are separate additions to the geometric observing workflows.

In the published v0.16.0 implementation, the analytical and SPK ephemerides
share an \texttt{EclipticMeanJ2000} type marker without an explicit transformation
between their native physical frame realizations. Type-level compatibility
therefore does not by itself establish physical realization equivalence. The
analytical-versus-SPK table reports backend-output differences including this
convention effect; resolving the physical-frame contract is separate from the
same-kernel SPK validation.

The original rotation grid uses prescribed UT1, while the compact time
extension checks measured IERS interpolation, EOP-backed UT1 and inserted-second
labels independently. Full-JD \texttt{f64} outputs have coarser temporal resolution than the native
split-second representation, so a zero residual in the table must be
interpreted at the full-JD resolution.
The EOP final-day extension and exact SPK endpoint rejection are identified
implementation boundaries. An exact-range guard in the EOP caller is a
practical mitigation; any upstream correction would require separate
verification. The evaluated v0.16.0 dependency graph is unchanged.
Complete polar-motion-corrected frame chains and
relativistic time transformations would extend the present coverage. More extensive dynamics and comparative
workflow benchmarks would complement the present arithmetic measurements.
These directions fit the architecture's separation of representations,
scientific models and application workflows. Public deposition of the underlying research data and independent technical
review remain relevant to publication and subsequent reproduction.
